\chapter{Metodologia badań}
\label{ch:metodologia}

\section{Scenariusze i metryki}
\label{sec:scenariusze}

\subsection{Scenariusze badawcze}

Badania podzielono na dwie grupy: weryfikację poprawności oraz pomiary wydajności
i~skalowania. Weryfikacja poprawności obejmowała sprawdzenie,
że wynik obliczeń jest niezależny od~liczby locale, na~które rozdzielono dziedzinę, pod~warunkiem
prawidłowej wymiany warstw brzegowych. Jako niezmiennik kontrolny wykorzystano sumę wartości pola
w~stanie końcowym, raportowaną przez program w~wierszu \texttt{final field: min=\ldots{}
max=\ldots{} sum=\ldots}. Ponieważ pole przyjmuje wartości w~przedziale $[1,2]$, a~jego
średnia jest bliska jedności, suma jest w~przybliżeniu równa liczbie komórek siatki.

Pomiary wydajności obejmowały cztery grupy scenariuszy, zestawione w~tabeli~\ref{tab:scenariusze}:
\begin{itemize}
  \item \textbf{skalowanie wątkowe w~obrębie węzła}: jeden locale (\texttt{-nl 1}),
        zmienna liczba wątków $1/2/4/8/16$, przeprowadzone dla kostki $1000^3$ na~węźle
        Pionier-lab8-1 o~20~wątkach.
  \item \textbf{skalowanie wątkowe przy 9 locale}: kostka $1000^3$, ustalona
        liczba locale, zmienna liczba wątków na~locale.
  \item \textbf{skalowanie po~liczbie węzłów}: kostka $1000^3$, $16$ wątków na~locale,
        liczba locale od~1 do~9.
  \item \textbf{skalowanie rozmiaru zadania}: 9 locale, $16$ wątków na~locale,
        krawędź kostki od~$125$ do~$1000$ (a~więc od~$125^3$ do~$1000^3$ komórek).
\end{itemize}

\begin{table}[htbp]
\centering
\caption{Grupy scenariuszy pomiarowych. Liczbę wątków w~pierwszej grupie oznaczono
symbolem $p$, liczbę locale symbolem $L$, a~krawędź kostki symbolem $N$.}
\label{tab:scenariusze}
\small
\begin{tabular}{@{}>{\raggedright\arraybackslash}p{2.8cm}>{\raggedright\arraybackslash}p{3.4cm}>{\raggedright\arraybackslash}p{4.0cm}>{\raggedright\arraybackslash}p{2.8cm}@{}}
\toprule
Scenariusz & Zmienna & Ustalone & Środowisko \\
\midrule
Wątki, $1000^3$       & $p\in\{1{,}2{,}4{,}8{,}16\}$ & $L{=}1$, $N{=}1000$, 10 kroków & węzeł Pionier-lab8-1 (20 wątków) \\
Wątki, 9 locale      & $p\in\{1{,}2{,}4{,}8{,}16\}$ & $L{=}9$, $N{=}1000$, 100 kroków   & klaster (9 węzłów) \\
Węzły                 & $L\in\{1{,}\ldots{,}9\}$ & $p{=}16$, $N{=}1000$, 100 kroków & klaster (9 węzłów) \\
Rozmiar kostki        & $N\in\{125{,}\ldots{,}1000\}$ & $p{=}16$, $L{=}9$, 100 kroków & klaster (9 węzłów) \\
\bottomrule
\end{tabular}
\end{table}

Skalowanie wątkowe w~obrębie węzła i~skalowanie po~węzłach rozdzielono świadomie:
to dwa różne mechanizmy równoległości, współdzielenie pamięci przez
wątki jednego locale oraz rozproszenie danych między locale połączone z~komunikacją sieciową.
Dopiero zestawienie obu perspektyw pozwala zidentyfikować czynnik ograniczający wydajność.

Górny pułap \num{16} wątków w~powyższych scenariuszach wynika z~architektury węzłów oraz
z~domyślnego zachowania warstwy wykonawczej Chapel. Procesor Intel~Core~i7-12700K jest
układem hybrydowym: ma 8~rdzeni wydajnościowych z~wielowątkowością współbieżną, co daje
\num{16} wątków sprzętowych, oraz 4~rdzenie energooszczędne bez wielowątkowości, łącznie
12~rdzeni fizycznych i~20 wątków logicznych. Skrypt uruchomieniowy żąda liczby wątków na~locale
równej wynikowi polecenia \texttt{nproc}, czyli 20. Warstwa qthreads, korzystająca z~biblioteki
hwloc, przy domyślnym ustawieniu \texttt{CHPL\_RT\_USE\_PU\_KIND=performance} wykorzystuje jednak
wyłącznie jednostki wykonawcze rdzeni wydajnościowych i~pomija rdzenie energooszczędne, co na~tych
węzłach daje \num{16} dostępnych wątków logicznych. Żądanie 20~wątków przekracza tę liczbę, więc
warstwa wykonawcza ogranicza ją do~\num{16} i~zgłasza komunikat \texttt{Reduced
numThreadsPerLocale=20 to 16 to prevent oversubscription}~\cite{chapel-docs}. Każdy węzeł pracuje
zatem efektywnie na~\num{16} wątkach i~tę wartość ustalono we~wszystkich przebiegach wielowęzłowych.

Z~tego samego powodu skalowanie wątkowe w~obrębie węzła prowadzono w~potęgach dwójki
$1/2/4/8/16$. Kolejne podwojenia dają równomiernie rozmieszczone punkty na~osi logarytmicznej,
a~górny punkt \num{16} pokrywa się z~pułapem narzuconym przez warstwę wykonawczą oraz z~liczbą
wątków na~locale używaną w~przebiegach wielowęzłowych, dzięki czemu obie serie pozostają
porównywalne. Nie uwzględniono 20~wątków, ponieważ rdzenie energooszczędne są domyślnie pomijane,
a~samo żądanie 20 i~tak jest redukowane do~\num{16}. Nie uwzględniono też 12~wątków
odpowiadających liczbie rdzeni fizycznych, ponieważ warstwa wykonawcza liczy jednostki wykonawcze
rdzeni wydajnościowych wraz z~wielowątkowością, a~nie same rdzenie fizyczne, a~ponadto 12 nie jest
potęgą dwójki. Jedyną zmienną niezależną w~tych scenariuszach była odpowiednio liczba węzłów lub
rozmiar dziedziny.

Wszystkie przebiegi wykorzystują ten sam zestaw parametrów numerycznych. Współczynnik
dyfuzji przyjęto bezwymiarowo $\alpha=\num{0.25}$, a~gorącą warstwę startową o~grubości
\texttt{hotThickness}=\num{2} płaszczyzn inicjalizowano wzdłuż osi~$y$ przy ścianie $y=0$.
Krok czasowy dobierany jest automatycznie z~warunku~\eqref{eq:cfl}, z~zastosowaniem
współczynnika bezpieczeństwa obniżającego krok o~\SI{2}{\percent} poniżej granicy stabilności:
\begin{equation}
  \Delta t = \num{0.98}\cdot\Delta t_{\max},
  \label{eq:dt}
\end{equation}
co dla~najbardziej niekorzystnego modu, o~naprzemiennym znaku między sąsiednimi
węzłami, daje współczynnik wzmocnienia~\eqref{eq:amp} równy $g=\num{-0.96}$. Ujemny
znak oznacza jedynie, że~amplituda tego modu zmienia znak w~każdym kroku, natomiast
warunek $\lvert g\rvert<1$ gwarantuje, że~jej moduł maleje. Błędy nie narastają,
lecz są tłumione, więc schemat pozostaje stabilny mimo błędów zaokrągleń.
Dla~największej badanej kostki $N=1000$ daje to $\Delta t\approx 6{,}55\cdot 10^{-7}$. Liczba kroków
czasowych $K$ zależy od~grupy scenariusza: wynosi \num{10} przy~skalowaniu wątkowym
na~pojedynczym węźle oraz \num{100} w~eksperymentach wielowęzłowych.

\subsection{Metryki}

Podstawową wielkością mierzoną jest \emph{czas wykonania} $T$, rozumiany jako czas trwania
całej pętli czasowej, mierzony zegarem \texttt{computeTimer} w~locale głównym (locale~0).
Ponieważ krok czasowy zawiera barierę synchronizacyjną wymiany halo, wszystkie locale kończą
krok w~tym samym czasie, więc pomiar z~locale głównego jest reprezentatywny dla całego
przebiegu. Czas ten nie obejmuje inicjalizacji, mierzonej osobno zegarem
\texttt{initTimer}.

Na podstawie czasu wykonania zdefiniowano standardowe miary skalowania silnego:
przyspieszenie $S(p)=T(1)/T(p)$ oraz wydajność równoległą $E(p)=S(p)/p$, zdefiniowane już
wzorami~\eqref{eq:speedup} i~\eqref{eq:efficiency}, gdzie $T(1)$ jest czasem wykonania konfiguracji
bazowej z~jednym wątkiem albo jednym locale.
W~roli $p$ dla skalowania wątkowego przyjęto \emph{efektywną} liczbę wątków,
czyli wartość \texttt{maxTaskPar} po~ewentualnym ograniczeniu przez warstwę wątkową, a~nie liczbę
żądaną. Dla oceny wykorzystania sprzętu niezależnie od~rozmiaru zadania posłużono się
\emph{przepustowością} wyrażoną w~miliardach aktualizacji komórek na~sekundę,
\begin{equation}
  \Pi = \frac{N^3 \cdot K}{T \cdot 10^{9}},
  \label{eq:throughput}
\end{equation}
gdzie $N^3 \cdot K$ jest łączną liczbą aktualizacji komórek (wielkość bezwymiarowa),
$T$ jest czasem wykonania w~sekundach, a~czynnik $10^{9}$ jest prefiksem giga
zamieniającym aktualizacje na~giga-aktualizacje. $\Pi$ wyrażone jest
w~giga-aktualizacjach komórek na~sekundę, co odpowiada konwencji metryki GUPS z~systemu
HPC Challenge~\cite{hpcchallenge}. Dla oszacowania roli
komunikacji wprowadzono \emph{udział komunikacji} w~czasie pojedynczego kroku,
\begin{equation}
  \texttt{comm\_frac} = \frac{t_{\text{fluff}}}{t_{\text{fluff}} + t_{\text{compute}}},
  \label{eq:commfrac}
\end{equation}
gdzie $t_{\text{fluff}}$ i~$t_{\text{compute}}$ są uśrednionymi po~krokach czasami wymiany halo
oraz obliczeń. Wielkość \texttt{comm\_frac} przyjmuje wartości bliskie jedności, gdy dominuje komunikacja,
oraz bliskie zeru, gdy dominują obliczenia. Uzupełniająco rejestrowano objętość wymiany halo
w~bajtach na~krok, obliczaną analitycznie przez procedurę \texttt{haloBytesPerStep} jako łączną
liczbę komórek warstw brzegowych wszystkich locale pomnożoną przez rozmiar elementu.

\subsection{Wybór backendu kompilatora}
\label{sec:backend}

Chapel generuje kod wynikowy jedną z~dwóch dróg: przez backend C (\texttt{CHPL\_LLVM=none})
albo przez backend LLVM (\texttt{CHPL\_LLVM=system}). Aby sprawdzić, czy wybór backendu wpływa
na~wydajność w~docelowym środowisku rozproszonym, ten sam program zbudowano w~obu wariantach
i~porównano na~dziewięciu węzłach klastra Pionier. Każdy wariant uruchomiono dziesięciokrotnie
dla kostki $100^3$, przez \num{100} kroków, z~warstwą komunikacyjną mpi-conduit. Backend LLVM
korzystał z~systemowej instalacji LLVM w~wersji 19.1.7. Wyniki zestawiono w~\cref{tab:llvm}.

\begin{table}[htbp]
\centering
\caption{Porównanie backendu kompilatora C (\texttt{CHPL\_LLVM=none}) i~LLVM
(\texttt{CHPL\_LLVM=system}, LLVM 19.1.7) dla kostki $100^3$ na~dziewięciu węzłach klastra
Pionier, \num{100} kroków, \num{10} powtórzeń. Kolumny mediana, średnia, min i~max dotyczą
czasu wykonania, a~kolumny updateFluff i~compute podają medianę czasu danej fazy na~krok.}
\label{tab:llvm}
\small
\begin{tabular}{@{}lrrrrrr@{}}
\toprule
Backend & mediana & średnia & min & max & updateFluff & compute \\
        & [s] & [s] & [s] & [s] & [s/krok] & [s/krok] \\
\midrule
C (\texttt{none}) & \num{20.015} & \num{19.995} & \num{19.883} & \num{20.096} & \num{0.1966} & \num{0.00185} \\
LLVM 19.1.7       & \num{19.910} & \num{19.908} & \num{19.842} & \num{19.962} & \num{0.1955} & \num{0.00185} \\
\bottomrule
\end{tabular}
\end{table}

Mediany czasu wykonania różnią się o~\SI{0.105}{\second}, czyli około \SI{0.5}{\percent}, co
mieści się w~granicach rozrzutu między powtórzeniami. Wyjaśnia tę różnicę rozkład czasu
pojedynczego kroku: wymiana halo zajmuje około \SI{99}{\percent} kroku, a~same obliczenia około
\SI{1}{\percent}, więc udział komunikacji \texttt{comm\_frac} według~\eqref{eq:commfrac} jest
bliski jedności. Backend kompilatora wpływa wyłącznie na~lokalny kod obliczeniowy, którego czas
jest w~obu wariantach praktycznie taki sam i~stanowi niewielką część kroku, dlatego nie przekłada
się na~całkowity czas wykonania.

Z~tego powodu we~wszystkich dalszych pomiarach używano backendu C (\texttt{CHPL\_LLVM=none}),
który w~komunikacyjnie ograniczonym środowisku wielowęzłowym nie powoduje mierzalnej straty
wydajności.

\section{Procedura i narzędzia pomiarowe}
\label{sec:procedura}

\subsection{Procedura pomiarowa}

Każdą konfigurację uruchamiano dziesięciokrotnie, a~jako wartość reprezentatywną przyjmowano \emph{medianę} czasu wykonania. Wybór mediany zamiast średniej arytmetycznej jest istotny:
pojedyncze przebiegi bywają zaburzone przez chwilowe zdarzenia systemowe, takie jak rywalizacja o~rdzenie,
przejściowe ograniczenie taktowania procesora czy aktywność innych procesów, które zawyżają
czas, lecz nie odzwierciedlają właściwości badanej konfiguracji. Mediana jest odporna na takie
wartości odstające, podczas gdy średnia zostaje przez nie zawyżona. Przykładem jest jedna
z~serii jednowęzłowych, w~której trzy z~dziesięciu przebiegów były wolniejsze o~około pół
sekundy (około \SI{5}{\percent} czasu pojedynczego przebiegu) wskutek przejściowego ograniczenia
taktowania. Skutkowało to podwyższonym odchyleniem
standardowym przy praktycznie niezmienionej medianie.

Przebiegi danej serii wykonywano sekwencyjnie, po~jednym naraz, aby uniknąć wzajemnej
rywalizacji o~pamięć i~pasmo. Serie pomiarowe wykonywano zautomatyzowanym skryptem
\path{bench.sh} opisanym w~\cref{sec:uruchamianie}. Wszystkie pomiary wydajności prowadzono na~środowisku Chapel~2.9
z~warstwą komunikacyjną w~wariancie mpi-conduit. Liczbę wątków na~locale ustawiano zmienną
środowiskową \texttt{CHPL\_RT\_NUM\_THREADS\_PER\_LOCALE}.

Konfigurację każdego przebiegu zapisywano bezpośrednio w~logu w~postaci linii
\texttt{[cfg]} opisanych w~\cref{sec:uruchamianie}: skrypt podaje żądaną liczbę wątków i~locale,
a~program efektywną liczbę wątków po~ograniczeniu przez warstwę qthreads. Na~węzłach Pionier
o~20~wątkach sprzętowych liczba efektywna była równa żądanej dla~wszystkich badanych
konfiguracji ($p\le 16$).


\subsection{Narzędzia pomiarowe}

Czasy mierzono wbudowanymi w~język stoperami (\texttt{stopwatch}), opartymi na~zegarze
monotonicznym wysokiej rozdzielczości o~praktycznej rozdzielczości pojedynczych
mikrosekund, osobno dla inicjalizacji, całej pętli czasowej oraz, w~każdym kroku,
dla wymiany halo, obliczeń i~zapisu. Objętość i~liczbę operacji komunikacyjnych GET, PUT, zdalnych i~atomowych rejestrowano za~pomocą modułu diagnostyki komunikacji GASNet,
włączanego przełącznikiem \texttt{commLog}. Wypisywany raz wiersz \texttt{[comm]} podaje
analityczną objętość wymiany halo na~krok. Zużycie pamięci procesu kontrolowano narzędziem
systemowym \texttt{/usr/bin/time -v}, odczytując maksymalny rezydentny zbiór stron
(ang.~\emph{Resident Set Size}, RSS), co pozwoliło ustalić
graniczne rozmiary kostek mieszczące się w~pamięci danego węzła.

Przetwarzanie zebranych logów oraz rysowanie wykresów zrealizowano skryptami w~języku Python
z~biblioteką \texttt{matplotlib}. Podstawową formą prezentacji jest wykres pudełkowy
rozkładu czasu wykonania z~naniesionymi punktami poszczególnych przebiegów, medianą
oraz średnią. Uzupełniają go wykresy metryk pochodnych: przyspieszenia, wydajności,
przepustowości i~udziału komunikacji. Taka forma pokazuje jednocześnie środek rozkładu
i~rozrzut, co jest istotne przy ocenie powtarzalności.

Do diagnozy zachowania sieci wykorzystano dedykowany skrypt sondujący, który równocześnie
obciąża łącze strumieniem \texttt{iperf3} w~trybie UDP o~przepływności zbliżonej do~granicy
łącza oraz monitoruje straty pakietów z~kilku niezależnych perspektyw: procentu utraconych
pakietów ICMP, przyrostów liczników odrzuceń na~lokalnym interfejsie sieciowym oraz, co
najważniejsze, przyrostów liczników błędów bufora odbiorczego UDP po~stronie odbiorcy
(\texttt{InErrors}, \texttt{RcvbufErrors}). Rosnące błędy bufora odbiorczego pod~obciążeniem są
bezpośrednim dowodem przeciążenia typu \emph{incast}, które w~wariancie AMUDP prowadzi do~awarii
typu \texttt{ECONGESTION}~\cite{incast}. Narzędzie to pozwoliło odróżnić przyczyny sieciowe od~obliczeniowych
i~dyskowych oraz uzasadnić wybór wariantu transportu opartego na~protokole TCP.

\subsection{Przykład wyznaczenia metryk}

Dla uporządkowania powyższych definicji poniżej przedstawiono sposób wyznaczenia metryk pochodnych na~jednym
scenariuszu. Rozważono skalowanie wątkowe kostki $1000^3$ na~węźle Pionier o~20~wątkach. Jednowątkowa
konfiguracja bazowa osiągnęła medianę czasu wykonania $T(1)=\SI{9.812}{\second}$, natomiast
optymalna konfiguracja ośmiowątkowa $T(8)=\SI{8.094}{\second}$. Zgodnie z~\eqref{eq:speedup}
przyspieszenie wynosi $S(8)=9{,}812/8{,}094\approx\num{1.212}$, a~wydajność równoległa według
\eqref{eq:efficiency}, $E(8)=1{,}212/8\approx\num{0.152}$. Przepustowość dla tej konfiguracji, przy
$K=10$ krokach, obliczono z~\eqref{eq:throughput}: $1000^3\cdot 10 = 10^{10}$ aktualizacji komórek
w~czasie $T(8)=\SI{8.094}{\second}$ daje $10^{10}/8{,}094 \approx \num{1.24}\cdot 10^{9}$ aktualizacji
na~sekundę, a~po~podzieleniu przez $10^{9}$ (prefiks giga) $\Pi \approx \num{1.24}$ miliarda
aktualizacji komórek na~sekundę.
Zabieg ten powtórzono dla każdej konfiguracji każdego scenariusza, dzięki czemu wszystkie wykresy
i~zestawienia w~\cref{ch:wyniki} opierają się na~jednolicie zdefiniowanych wielkościach.

\subsection{Kontrola warunków i~zagrożenia dla trafności}

Wiarygodność wniosków zależy od~kontroli czynników zakłócających. W~toku badań zidentyfikowano
i~ograniczono kilka takich czynników.

Zmienność środowiska pochodzi ze~współdzielenia maszyny z~innymi procesami oraz dynamicznego
zarządzania taktowaniem procesora, co wprowadza zmienność czasu wykonania. Ograniczono ją przez
sekwencyjne uruchamianie przebiegów, dziesięciokrotne powtórzenia oraz stosowanie mediany. Tam,
gdzie mimo to pojawiły się wartości odstające, czyli pojedyncze wolniejsze przebiegi wskutek przejściowego
ograniczenia taktowania, fakt ten odnotowano jawnie i~potraktowano jako artefakt pomiaru, a~nie
własność konfiguracji.

Nierówny podział dziedziny pojawia się, gdy liczba locale nie dzieli równo boku kostki $1000^3$.
Powoduje to nierównomierne obciążenie i~zaburza gładkość krzywej skalowania. Nie jest to
błąd pomiaru, lecz rzeczywista własność dekompozycji. Uwzględniono ją w~interpretacji
niemonotonicznego przebiegu skalowania po~węzłach.

Trafność zewnętrzna wyników jest ograniczona: uzyskano je na~konkretnym, niewielkim klastrze z~siecią
\SI{1}{\giga\bit\per\second}. Wartości bezwzględne nie przenoszą się wprost na~inne środowiska. Wnioski mają
jednak charakter ogólny i~wynikają z~właściwości samych obliczeń stencilowych:
obliczenia te są ograniczone przepustowością pamięci oraz komunikacją sieciową, a~większe
zadania lepiej wykorzystują klaster~\cite{roofline,datta-stencil}.
